\lecture{21}{Двойственные и седловые методы}

\sourcevideo
    {https://www.youtube.com/playlist?list=PLOzRYVm0a65fhKDS8cYIC7YCuVGJ4FOJx}
    {Week 6: Lecture 21: Constrained Optimization Problem}

Задача с ограничениями может быть преобразована в двойственную задачу
по множителям Лагранжа или в задачу поиска седловой точки. Для
линейных равенств двойственная функция выражается через сопряжённую
функцию Фенхеля, а её градиент совпадает с невязкой ограничения.

\section{Функция Лагранжа и двойственная задача}

Рассмотрим задачу
\[
    \min_{x\in\mathbb R^n} f_0(x)
\]
при условиях
\[
    f_i(x)\le0,
    \qquad
    i=1,\ldots,m,
\]
\[
    h_j(x)=0,
    \qquad
    j=1,\ldots,r.
\]
Переменная \(x\) называется прямой. Для неравенств вводятся множители
\(\lambda_i\ge0\), а для равенств --- свободные множители
\(\nu_j\in\mathbb R\).

\begin{definition}[Функция Лагранжа]
Функцией Лагранжа называется
\[
    \mathcal L(x,\lambda,\nu)
    =
    f_0(x)
    +
    \sum_{i=1}^{m}\lambda_i f_i(x)
    +
    \sum_{j=1}^{r}\nu_j h_j(x).
\]
Переменные \(\lambda\) и \(\nu\) называются двойственными.
\end{definition}

\begin{definition}[Двойственная функция]
Двойственная функция определяется формулой
\[
    g(\lambda,\nu)
    =
    \inf_{x\in\mathbb R^n}
    \mathcal L(x,\lambda,\nu).
\]
\end{definition}

При каждом фиксированном \(x\) функция
\[
    (\lambda,\nu)
    \longmapsto
    \mathcal L(x,\lambda,\nu)
\]
аффинна. Поэтому её поточечный инфимум \(g\) всегда вогнут, даже если
исходная задача невыпукла. Двойственная задача имеет вид
\[
    \max_{\lambda\ge0,\,\nu}g(\lambda,\nu).
\]

Для любой допустимой пары \((\lambda,\nu)\) выполняется слабая
двойственность
\[
    g(\lambda,\nu)\le p^\star,
\]
где \(p^\star\) --- оптимальное значение прямой задачи. Если
\[
    d^\star=p^\star,
\]
двойственный разрыв отсутствует. Для выпуклых задач это обычно
обеспечивается условием регулярности, например условием Слейтера.

\section{Сопряжённая функция Фенхеля}

\begin{definition}[Сопряжённая функция]
Для функции
\[
    f\colon\mathbb R^n
    \to
    \mathbb R\cup\{+\infty\}
\]
сопряжённая функция Фенхеля определяется как
\[
    f^\ast(y)
    =
    \sup_{x\in\mathbb R^n}
    \left\{
        \langle y,x\rangle-f(x)
    \right\}.
\]
\end{definition}

Для каждого \(x\) выражение
\[
    \langle y,x\rangle-f(x)
\]
аффинно по \(y\), поэтому \(f^\ast\) всегда выпукла. Геометрически
\(f^\ast(y)\) является максимальным разрывом между линейной функцией
\[
    x\longmapsto\langle y,x\rangle
\]
и функцией \(f\).

Для замкнутых собственных выпуклых функций сильная выпуклость и
гладкость переходят к сопряжённой функции с обратными константами.

\begin{theorem}
Если \(f\) является \(\mu\)-сильно выпуклой, то \(f^\ast\)
дифференцируема и
\[
    \left\lVert
        \nabla f^\ast(y_1)
        -
        \nabla f^\ast(y_2)
    \right\rVert
    \le
    \frac1\mu
    \lVert y_1-y_2\rVert.
\]
Кроме того,
\[
    \nabla f^\ast(y)
    =
    \operatorname*{argmax}_{x}
    \left\{
        \langle y,x\rangle-f(x)
    \right\}.
\]
\end{theorem}

\begin{theorem}
Если \(f\) выпукла и \(L\)-гладка, то \(f^\ast\) является
\(1/L\)-сильно выпуклой на своей области определения.
\end{theorem}

Таким образом, сильная выпуклость исходной функции обеспечивает
гладкость сопряжённой, а гладкость исходной функции --- сильную
выпуклость сопряжённой.

\section{Линейные равенства и двойственный подъём}

Рассмотрим задачу
\[
    \min_{x\in\mathbb R^n}f(x)
    \qquad
    \text{при условии}
    \qquad
    Ax=b,
\]
где
\[
    A\in\mathbb R^{m\times n},
    \qquad
    b\in\mathbb R^m.
\]
Её функция Лагранжа равна
\[
    \mathcal L(x,\nu)
    =
    f(x)+\nu^{\mathsf T}(Ax-b).
\]

Двойственная функция вычисляется следующим образом:
\[
\begin{aligned}
    g(\nu)
    &=
    \inf_x
    \left\{
        f(x)+\nu^{\mathsf T}(Ax-b)
    \right\}
    \\
    &=
    -b^{\mathsf T}\nu
    +
    \inf_x
    \left\{
        f(x)+\langle A^{\mathsf T}\nu,x\rangle
    \right\}
    \\
    &=
    -f^\ast(-A^{\mathsf T}\nu)
    -
    b^{\mathsf T}\nu.
\end{aligned}
\]
Следовательно, двойственная задача не имеет ограничений по \(\nu\):
\[
    \max_{\nu\in\mathbb R^m}
    \left\{
        -f^\ast(-A^{\mathsf T}\nu)
        -
        b^{\mathsf T}\nu
    \right\}.
\]

Пусть \(f\) сильно выпукла. Тогда
\[
    x(\nu)
    =
    \operatorname*{argmin}_x
    \mathcal L(x,\nu)
\]
единственен и
\[
    x(\nu)
    =
    \nabla f^\ast(-A^{\mathsf T}\nu).
\]
По теореме Данскина
\[
    \nabla g(\nu)=Ax(\nu)-b.
\]
Таким образом, градиент двойственной функции совпадает с невязкой
прямого ограничения.

Обычный двойственный подъём имеет вид
\[
    \nu^{k+1}
    =
    \nu^k
    +
    \eta
    \bigl(
        Ax(\nu^k)-b
    \bigr).
\]
На каждой итерации сначала решается внутренняя задача по \(x\), затем
множитель изменяется в направлении нарушения ограничения.

Пусть вогнутая функция \(g\) удовлетворяет PL-условию для
максимизации:
\[
    \frac12
    \lVert\nabla g(\nu)\rVert^2
    \ge
    \mu_d
    \bigl(
        g^\star-g(\nu)
    \bigr).
\]
Для
\[
    p>2,
    \qquad
    1<q<2
\]
положим
\[
    \rho_p=\frac{p-2}{p-1},
    \qquad
    \rho_q=\frac{q-2}{q-1}.
\]
Фиксированно-временной двойственный поток записывается как
\[
\begin{aligned}
    \dot\nu
    ={}&
    c_1
    \frac{\nabla g(\nu)}
    {\lVert\nabla g(\nu)\rVert^{\rho_p}}
    \\
    &+
    c_2
    \frac{\nabla g(\nu)}
    {\lVert\nabla g(\nu)\rVert^{\rho_q}},
\end{aligned}
\]
где \(c_1,c_2>0\). Для
\[
    V(\nu)=g^\star-g(\nu)
\]
имеем
\[
    \dot V
    \le
    -c_1(2\mu_d)^{\alpha_1}V^{\alpha_1}
    -
    c_2(2\mu_d)^{\alpha_2}V^{\alpha_2},
\]
где
\[
    \alpha_1
    =
    \frac{p}{2(p-1)}
    <1,
    \qquad
    \alpha_2
    =
    \frac{q}{2(q-1)}
    >1.
\]
По критерию фиксированного времени двойственный оптимум достигается за
время, равномерно ограниченное по начальному множителю.

\section{Квадратичная задача и система KKT}

Рассмотрим
\[
    \min_x
    \left\{
        \frac12x^{\mathsf T}Qx+c^{\mathsf T}x
    \right\}
    \qquad
    \text{при условии}
    \qquad
    Ax=b,
\]
где
\[
    Q\succ0.
\]
Для
\[
    f(x)
    =
    \frac12x^{\mathsf T}Qx+c^{\mathsf T}x
\]
сопряжённая функция равна
\[
    f^\ast(y)
    =
    \frac12
    (y-c)^{\mathsf T}
    Q^{-1}
    (y-c).
\]
Действительно, условие оптимальности во внутренней задаче даёт
\[
    y-Qx-c=0,
    \qquad
    x=Q^{-1}(y-c).
\]

Поэтому двойственная функция имеет вид
\[
\begin{aligned}
    g(\nu)
    ={}&
    -\frac12
    \bigl(
        c+A^{\mathsf T}\nu
    \bigr)^{\mathsf T}
    Q^{-1}
    \bigl(
        c+A^{\mathsf T}\nu
    \bigr)
    \\
    &-
    b^{\mathsf T}\nu,
\end{aligned}
\]
а минимизатор функции Лагранжа равен
\[
    x(\nu)
    =
    -Q^{-1}
    \bigl(
        c+A^{\mathsf T}\nu
    \bigr).
\]
Следовательно,
\[
    \nabla g(\nu)
    =
    -AQ^{-1}
    \bigl(
        c+A^{\mathsf T}\nu
    \bigr)
    -
    b
\]
и
\[
    \nabla^2g(\nu)
    =
    -AQ^{-1}A^{\mathsf T}
    \preceq0.
\]
Если \(A\) имеет полный строчный ранг, то
\[
    AQ^{-1}A^{\mathsf T}\succ0,
\]
поэтому \(g\) строго вогнута.

Условия оптимальности равны
\[
    Qx^\star+c+A^{\mathsf T}\nu^\star=0,
    \qquad
    Ax^\star=b.
\]
Они образуют систему KKT
\[
\begin{pmatrix}
    Q&A^{\mathsf T}\\
    A&0
\end{pmatrix}
\begin{pmatrix}
    x^\star\\
    \nu^\star
\end{pmatrix}
=
\begin{pmatrix}
    -c\\
    b
\end{pmatrix}.
\]
При \(Q\succ0\) и полном строчном ранге \(A\) пара
\[
    (x^\star,\nu^\star)
\]
единственна.

\section{Седловая динамика и ньютоновский поток}

Для задачи
\[
    \min_x\max_z\Phi(x,z)
\]
пара \((x^\star,z^\star)\) называется седловой точкой, если
\[
    \Phi(x^\star,z)
    \le
    \Phi(x^\star,z^\star)
    \le
    \Phi(x,z^\star)
\]
для всех допустимых \(x,z\). Для линейного ограничения функция
Лагранжа задаёт именно такую постановку:
\[
    \min_x\max_\nu\mathcal L(x,\nu).
\]

Естественная динамика спуска--подъёма имеет вид
\[
    \dot x=-\nabla_x\Phi(x,z),
    \qquad
    \dot z=\nabla_z\Phi(x,z).
\]
Введём
\[
    w=
    \begin{pmatrix}
        x\\
        z
    \end{pmatrix},
    \qquad
    F(w)=
    \begin{pmatrix}
        \nabla_x\Phi(x,z)\\
        -\nabla_z\Phi(x,z)
    \end{pmatrix}.
\]
Тогда седловая точка удовлетворяет
\[
    F(w^\star)=0,
\]
а спуск--подъём записывается как
\[
    \dot w=-F(w).
\]

Якобиан седлового оператора равен
\[
    J_F(w)
    =
    \begin{pmatrix}
        \nabla^2_{xx}\Phi
        &
        \nabla^2_{xz}\Phi
        \\
        -\nabla^2_{zx}\Phi
        &
        -\nabla^2_{zz}\Phi
    \end{pmatrix}.
\]
Если \(\Phi\) строго выпукла по \(x\), строго вогнута по \(z\), а
смешанные производные согласованы, то
\[
    \frac{
        J_F(w)+J_F(w)^{\mathsf T}
    }2
    =
    \begin{pmatrix}
        \nabla^2_{xx}\Phi&0\\
        0&-\nabla^2_{zz}\Phi
    \end{pmatrix}
    \succ0.
\]
Следовательно, \(J_F(w)\) невырожден.

Ньютоновский седловой поток имеет вид
\[
    \dot w=-J_F(w)^{-1}F(w).
\]
Для
\[
    V(w)=\frac12\lVert F(w)\rVert^2
\]
получаем
\[
\begin{aligned}
    \dot V
    &=
    F(w)^{\mathsf T}J_F(w)\dot w
    \\
    &=
    -\lVert F(w)\rVert^2
    =
    -2V.
\end{aligned}
\]
Поэтому невязка седловых условий затухает экспоненциально.

Фиксированно-временная модификация задаётся уравнением
\[
\begin{aligned}
    \dot w
    ={}&
    -c_1J_F(w)^{-1}
    \frac{F(w)}
    {\lVert F(w)\rVert^{\rho_p}}
    \\
    &-
    c_2J_F(w)^{-1}
    \frac{F(w)}
    {\lVert F(w)\rVert^{\rho_q}},
\end{aligned}
\]
а при \(F(w)=0\) полагается \(\dot w=0\). Тогда
\[
    \dot V
    =
    -c_1 2^{\alpha_1}V^{\alpha_1}
    -
    c_2 2^{\alpha_2}V^{\alpha_2}
\]
и
\[
    T_{\max}
    =
    \frac1{
        c_1 2^{\alpha_1}(1-\alpha_1)
    }
    +
    \frac1{
        c_2 2^{\alpha_2}(\alpha_2-1)
    }.
\]
Если корень \(F\) единственен, то
\[
    w(t)=w^\star
\]
после времени \(T_{\max}\).

\section{Предпосылки и переход к дополненной функции}

Строгая выпуклость по \(x\) и строгая вогнутость по \(z\)
обеспечивают не более одной седловой точки, но не всегда гарантируют
её существование на неограниченной области. Дополнительно могут
потребоваться компактность, коэрцитивность, антикоэрцитивность или
условия сильной двойственности. Для ньютоновской динамики также
предполагаются существование траектории и обратимость \(J_F(w)\)
вдоль неё.

Двойственный метод требует на каждой итерации вычислять
\[
    x(\nu)
    =
    \operatorname*{argmin}_x
    \mathcal L(x,\nu)
\]
или, что эквивалентно при сильной выпуклости, значение
\[
    \nabla f^\ast(-A^{\mathsf T}\nu).
\]
В общем случае это самостоятельная задача оптимизации.

Ньютоновский седловой поток требует решать линейную систему
\[
    J_F(w)d=-F(w).
\]
Матрицу явно не обращают. После дискретизации фиксированно-временные
гарантии требуют отдельного анализа шага и устойчивости.

Обычная функция Лагранжа
\[
    \mathcal L(x,\nu)
    =
    f(x)+\nu^{\mathsf T}(Ax-b)
\]
может иметь недостаточную кривизну по прямой переменной. В следующей
лекции к ней добавляется квадратичный штраф:
\[
    \mathcal L_\rho(x,\nu)
    =
    f(x)
    +
    \nu^{\mathsf T}(Ax-b)
    +
    \frac\rho2
    \lVert Ax-b\rVert^2.
\]
Эта дополненная функция Лагранжа лежит в основе метода множителей и
ADMM.
